% Source: https://artofproblemsolving.com/community/c3h3158474_2016_lmt_spring_guts_round_p13p24__lexington_mathematical_tournament
% Stable identity: aops-topic-3158474-q24
The remaining angles are $C=60^\circ,D=30^\circ$. Put $D=(0,0),C=(8,0)$ and let $a=AB,b=BC,d=DA$. The rays $DA,CB$ make angles $30^\circ,120^\circ$ with the positive horizontal axis, while $AB$ points at $-30^\circ$. Comparing vertical and horizontal components yields $d=a+\sqrt3b$ and $\sqrt3(a+d)+b=16$. Therefore $b=4-\dfrac{\sqrt3a}{2}$ and $d=4\sqrt3-\dfrac{a}{2}$. Splitting along $BD$, the area is $\dfrac{1}{2}ad\sin120^\circ+\dfrac{1}{2}(8)b\sin60^\circ=\dfrac{\sqrt3}{4}(ad+8b)$. Since $ad+8b=32-\dfrac{a^2}{2}$, the area condition gives $a^2=16$, so $a=4$. Then $b=4-2\sqrt3,d=4\sqrt3-2$, both positive. The perimeter is $\boxed{14+2\sqrt3}$.
\begin{asy}
usepackage("amsmath");
size(7cm);
real a=4,b=4-2sqrt(3),d=4sqrt(3)-2;pair D=(0,0),C=(8,0),A=(sqrt(3)*d/2,d/2),B=(8-b/2,sqrt(3)*b/2);draw(A--B--C--D--cycle);draw(B--D,dashed);label("$A$",A,N);label("$B$",B,NE);label("$C$",C,SE);label("$D$",D,SW);label("$8$",(4,0),S);label("$a$",(A+B)/2,N);label("$d$",(A+D)/2,W);label("$b$",(B+C)/2,E);
\end{asy}
